\section{Verwendung von KI-Werkzeugen}

Diese Dokumentenserie ist unter Verwendung KI-gestützter Werkzeuge entstanden.
Der Hinweis steht am Anfang, nicht im Kleingedruckten, und gehört als eigener
Abschnitt in jedes Dokument der Serie.

Die Arbeitsteilung ist eindeutig. Die Fragestellungen, die theoretischen
Setzungen, die geometrischen Ideen und alle inhaltlichen Entscheidungen stammen
vom Autor. Die KI-Werkzeuge sind Werkzeuge im wörtlichen Sinn: sie formulieren
aus, rechnen nach, prüfen auf Widerspruch, erstellen Prüfskripte, finden
Inkonsistenzen zwischen Dokumenten und übersetzen zwischen Darstellungen. Sie
treffen keine inhaltlichen Entscheidungen und erzeugen keine Thesen.

Diese Trennung ist nicht bloß eine Zuschreibung, sondern methodisch tragend.
Ein Sprachmodell erzeugt plausibel klingenden Text auch dort, wo keine Substanz
ist -- genau die Fehlerart, gegen die die Schichtmarkierung des nächsten
Abschnitts gerichtet ist. Deshalb gilt in dieser Serie durchgehend: jede
rechnerische Aussage hat ein hinterlegtes Prüfskript, das unabhängig
nachgerechnet werden kann; jede Zahl hat einen Toleranzmaßstab; jede
Herleitungsstufe hat eine Schichtmarke. Was diese Prüfung nicht besteht, steht
nicht als Ergebnis im Text.

Die Verantwortung für Inhalt und Fehler liegt vollständig beim Autor.


\section{Was nicht Teil der A-Serie ist}

Zwei Bereiche sind ausdrücklich ausgeschlossen.

\textbf{Der kosmologische Sektor.} Die A-Serie enthält ihn nicht. Die
entsprechenden Arbeiten des Altbestands bleiben bestehen und behalten ihre
Nummern, werden aber nicht in die neue Struktur übernommen. Der Grund ist der
Anspruchsunterschied: der geometrische Kern rechnet aus $\xi$, der
kosmologische Sektor arbeitet durchgehend mit importierten Skalen und
deklarierten Kalibrierungen. Beides in einer Serie zu führen, verwischt genau
die Trennung, um derentwillen die Serie neu aufgesetzt wird.

\textbf{Prognosen über die Grundgrößen hinaus.} Als Vorhersage geführt werden
allein die Größen, die unmittelbar aus dem Apparat folgen und an denen die
Theorie scheitern kann: die Tau-Masse, die $\xi$-Ordnung der
CHSH-Abweichung, die fraktale Dimension $D_f$ und die Grundlänge $L_0$.
Alles Weitere, was in der älteren Serie stellenweise als Prognose formuliert
war, erscheint hier als Strukturaussage oder gar nicht.

\section{Warum eine neue Serie}

Die bisherige Dokumentenserie der FFGFT ist historisch gewachsen. Ihre Nummern
folgen der Entstehungszeit, nicht der Sachlage. Über etwa drei Jahre sind so
mehr als zweihundert deutschsprachige Dokumente entstanden, die einzeln
nachvollziehbar sind, in ihrer Gesamtheit aber vier Eigenschaften angenommen
haben, die dem Ganzen schaden.

\textbf{Erstens Referenzverfall.} Ein erheblicher Teil der frühen Dokumente
wird von den späteren kaum noch zitiert. Der Inhalt ist nicht falsch geworden,
aber er ist an anderer Stelle vollständiger und präziser dargestellt. Die alte
Nummer bleibt dennoch stehen und wird weiter gefunden.

\textbf{Zweitens liegen die Korrekturen außerhalb der Texte.} Ein eigenes
Korrekturregister berichtigt die Dokumente von außen, statt sie zu ändern; es
führt Korrekturen, Präzisierungen und Revisionen als datierte Einträge. Wer ein
einzelnes altes Dokument liest, liest daher den unkorrigierten Stand. Die
Arbeitsregel, ältere Dokumente nicht zu verändern, war methodisch richtig --
sie hält die Entwicklungsgeschichte ehrlich fest --, aber sie verlagert die
Bringschuld auf die Leserin.

\textbf{Drittens Mehrfachbelegung.} Der Massensektor steht in mindestens fünf
Dokumenten, der Zeitbegriff in sieben, die Informationsfrage in sechs. Jede
dieser Darstellungen hat ihren Anlass, aber es gibt keinen Ort, an dem steht:
\emph{hier} ist der Stand.

\textbf{Viertens sind die Schichten typografisch unsichtbar.} Eine aus
$\xi$ hergeleitete Aussage, eine algebraisch bewiesene Übersetzung und eine
Plausibilitätsskizze stehen im selben Schriftschnitt nebeneinander. Das ist die
gefährlichste der vier Eigenschaften, weil sie den Anspruch der Theorie nach
oben verzerrt.

\textbf{[SETZUNG]} Die A-Serie beantwortet alle vier Punkte mit derselben Maßnahme: Sachordnung
statt Zeitordnung, Korrekturen eingearbeitet statt angehängt, ein Thema an
einem Ort, Schichtstatus pro Aussage ausgewiesen.

\section{Einstiegslektüre vor der A-Serie}

Wer die A-Serie liest, trifft ab A020 sofort auf Setzungen, Schichtmarkierungen
und geometrische Strukturen. Zwei Dokumente im Altbestand bereiten darauf vor --
auf verschiedene Weise.

\textbf{,,Die Reise'' (Altbestand, populärwissenschaftlich)} beschreibt die
T0-Theorie in Alltagssprache -- mit Analogien, ohne Formeln, ohne
Physik-Vorkenntnisse. Wer es vor A010 liest, kommt mit einem Bild im Kopf
an: warum eine einzige Zahl $\xi$ genügen könnte, was Zeit-Masse-Dualität
bedeutet, was an diesem Ansatz anders ist als an den etablierten Theorien.

\textbf{,,FFGFT in Alltagssprache'' (Altbestand, narrativ-technisch)} ist die
narrative Einführung in den Lagrangian-Kern. Es übersetzt die zentralen
Lagrangian-Formulierungen in zugängliche Sprache -- mit Bildern aus der
Akustik (schwingende Saite, Resonator, pythagoreisches Komma), die in der
FFGFT selbst als strukturelle Parallelen verwendet werden, nicht als beliebige
Metaphern. Jede Aussage ist an den formalen Dokumenten verankert. Wer verstehen
will, wie die Theorie rechnet, bevor er die Rechnungen sieht, liest dieses
Dokument.

Beide Dokumente sind nicht Teil der A-Serie und werden nicht durch sie ersetzt.
Sie sind eigenständig und bleiben es.

\section{Was die A-Serie nicht ist}

\textbf{[SETZUNG]} Sie ist \emph{keine} Revision im Sinne einer Neubewertung. Es wird nichts
zurückgenommen, was nicht schon im Korrekturregister zurückgenommen wurde. Der
Inhalt ist derselbe; die Ordnung ist neu.

Sie ersetzt die alten Dokumente auch nicht physisch. Der Altbestand bleibt
vollständig erhalten, unverändert, unter seinen alten Nummern. Er ist die
Entwicklungsgeschichte und bleibt zitierbar. Die A-Serie ist die Fassung, die
man liest, wenn man wissen will, was die Theorie behauptet -- nicht, wie sie
dazu gekommen ist.

Und sie ist keine Kürzung. Wo ein alter Text ausführlicher ist als der
entsprechende A-Abschnitt, verweist dieser darauf. Verloren gehen soll nur die
Redundanz, nicht die Substanz.

\section{Die drei Schichten}

Die FFGFT hat seit jeher drei methodische Schichten. Neu ist allein, dass sie
in der A-Serie markiert werden -- und zwar pro Aussage, nicht pro Dokument.

\begin{description}
  \item[{[K]} Kern.] Aus der Geometrie und dem einen Parameter $\xi$
    hergeleitet. Beispiel: die Koide-Relation $Q = 2/3$ als Operatorbedingung
    $r = \sqrt{2}$; die Rekursionsnichtschließung; $D_f = 3-\xi$.
  \item[{[B]} Brücke.] Algebraisch bewiesene Übersetzung zwischen zwei
    Darstellungen. Beispiel: die Bijektion $\mathcal{H} = L^2(T^4)\otimes
    \mathbb{C}^3$; die DFT-Diagonalisierung des $Z_3$-Zirkulanten. Eine Brücke
    trägt alles, was auf der einen Seite steht, und beweist nichts, was dort
    nicht schon stand.
  \item[{[S]} Skizze.] Plausibilität, Analogie, Kandidat. Beispiel: der
    Neutrinosektor; $N_0 \approx 38{,}6$; der Feinwert der Brechungsstärke $\delta$.
\end{description}

Dazu kommt eine vierte Markierung, die keine Schicht ist, sondern eine
Buchungsart:

\begin{description}
  \item[{[SETZUNG]}] Deklariert, nicht hergeleitet. Beispiel: $\tilde T\cdot m
    = 1$; die Wahl des $T^4$; der Zahlenwert $\xi = 4/30000$; der
    Referenzpunkt $m_\mu/m_e$.
\end{description}

\textbf{[B]} Der Sinn der Markierung ist nicht Bescheidenheit, sondern Prüfbarkeit. Eine
Theorie, die ihre Setzungen benennt, kann daran gemessen werden, wie wenige sie
braucht. Eine Theorie, die sie verschweigt, kann das nicht.

\textbf{[SETZUNG]} Die Marker [K] und [B] sind interne Zustandsangaben des
Anspruchs, keine externen Zertifizierungen. [K] besagt, dass eine Aussage
innerhalb des deklarierten Kerns folgt; [B] besagt, dass eine Brücke unter
ihren genannten Bedingungen algebraisch etabliert ist. Keiner der beiden
Marker begründet für sich empirische Wahrheit, externe physikalische
Gültigkeit oder eine eindeutige Ontologie.

\section{Der Toleranzmaßstab}

Jede Zahlenaussage der A-Serie nennt den Maßstab, gegen den sie geprüft wird.
Das ist eine Konsequenz aus einer Erfahrung, die älter ist als die Theorie: ein
Messinstrument ist stets weniger leistungsfähig als die Wirklichkeit, die es
misst. Punktgenaue Landung auf der letzten Ziffer ist deshalb weder erreichbar
noch das Ziel.

\textbf{[SETZUNG]} Vollständigkeit heißt in der A-Serie: \emph{es bleibt kein Rest, der über die
Toleranzen hinausragt und nur durch einen freien Parameter geschlossen werden
kann.} Ein Rest innerhalb der Bestimmungstoleranz ist kein Mangel, sondern der
Normalfall. Ein Rest, der die Toleranz übersteigt, ist eine Aussage -- und wird
als solche gebucht, nicht weggerundet.

Praktisch heißt das: bei jeder Abweichung stehen drei Angaben zusammen -- der
Wert, die Messgenauigkeit der Vergleichsgröße und die Bestimmungstoleranz der
Theoriegröße. Erst das Verhältnis dieser drei entscheidet, ob eine Abweichung
etwas bedeutet.

\section{Zitierregel der Serie}

\textbf{[SETZUNG]} Die A-Serie führt keine Literaturverzeichnisse. Die Regel:

\textbf{Standardphysik ohne Quelle.} Lehrbuchwissen -- Konstanten, Standardformeln,
etablierte Ergebnisse (Dirac-Gleichung, Casimir-Formel, CHSH-Schranke,
Koide-Relation, Landauer-Prinzip) -- wird ohne Zitat verwendet. Namensnennungen
(Bell, Tsirelson, Weyl) sind Konzeptbezeichnungen, keine Literaturverweise.

\textbf{Spezifische externe Messungen mit Kurzquelle im Text.} Wo ein konkretes,
datiertes Messergebnis Dritter mit Fehlerbalken behauptet wird, steht die
Quelle in Klammern direkt an der Zahl.

\textbf{Eigene Messungen mit Verweis auf das Messprotokoll.} Hardware-Läufe
und eigene Messkampagnen verweisen auf Rohdaten und Auswerteskripte im Korpus,
damit die Behauptung unabhängig prüfbar ist.

\section{Aufbau und Nummerierung}

Die Serie hat vier Blöcke.

\textbf{Block 0 (A010--A090): Grundlage.} Setzungen, Geometrie, fraktale
Korrektur, Rekursion, Zeit, Hilbertraum, Feldtheorie, Einheiten, Projektionskette.
Wer diesen Block gelesen hat, kennt den Apparat vollständig -- ohne ein einziges
Ergebnis. Aktuelle Dokumente: A010, A015, A020, A030, A040, A050, A060, A070,
A075, A080, A085, A090.

\textbf{Block 1 (A100--A190): Sektoren.} Leptonen, Koide-Relation,
$\theta = 2/9$, Feinstrukturkonstante, g-2, Gravitation, Gravitationskonstante,
Quarks und Neutrinos, Bell und Hardware, Information, Standardmodell. Hier
stehen die Ergebnisse, jedes auf den Apparat aus Block~0 zurückgeführt.
Aktuelle Dokumente: A100, A105, A110, A120, A130, A135, A138, A140, A142,
A145, A150, A160, A165, A180, A190.

\textbf{Block 2 (A200--A250): Methodik und Bilanz.} Die Unterscheidung von
Gesetz, Norm und Setzung; die Prüfgrößen; der Toleranzmaßstab; die offenen
Punkte; die Abgrenzung zu externen Rahmen; das Register.
Aktuelle Dokumente: A200, A210, A220, A230, A240, A250.

\textbf{Block 3 (A260--): Erweiterungen.} Themen ohne Rückverweis aus den
früheren Blöcken -- eigenständige Einheiten, die den Kern ergänzen, aber nicht
in den Hauptfluss eingebettet sind. Aktuelle Dokumente: A260 (Casimir-Effekt),
A261 (Skalenhierarchie), A262 ($c$-Konvention und $E=m$), A263 (Dirac
vollständig), A264 (Asymmetrie und CP), A265 (Statische Verschiebung und SI-Anker),
A266 (Einheitenprüfung und SI-Rückrechnung), A267 ($\alpha=1$ als Stipulation).

\textbf{Methodenstatus Block~3.} A260--A264 (Casimir, Skalenhierarchie,
$c$-Konvention, Dirac, Asymmetrie) und A265 (Spektralverschiebung und SI-Anker) sind
Physik-Ergebnisdokuemente ohne Rückverweis aus Block~0/1. A266 (Einheitenprüfung)
und A267 ($\alpha=1$-Konvention) sind Konventions- und Methodendokumente -- sie
enthalten Rechenwerkzeuge, keine neuen Theorieergebnisse. Das ändert nichts an
ihrem Platz in Block~3, aber der Leser sollte wissen: der Anspruch ist verschieden.

\textbf{[SETZUNG]} Die Nummern steigen innerhalb der Blöcke in Zehnerschritten.
Zwischennummern (A015, A075, A138, A142 u.\,a.) sind dann vergeben, wenn ein
späteres Dokument inhaltlich auf das frühere verweist und die Einordnung das
verlangt. Ein Dokument ohne Rückverweis aus dem bestehenden Block kommt an das
Ende -- in Block~3, nicht als Zwischennummer. Das ist die Regel, an der sich
die Entscheidung orientiert: nicht die Thematik, sondern der Verweis.

\section{Was von wo kommt}

Jedes A-Dokument nennt in seinem Schlussabschnitt die alten Nummern, deren
Inhalt es aufnimmt. Das ist keine Literaturliste, sondern eine
Ersetzungserklärung: wer das A-Dokument gelesen hat, muss die genannten alten
Dokumente nicht mehr lesen, um den Stand zu kennen.

Die umgekehrte Richtung -- von jeder alten Nummer zu ihrem A-Nachfolger --
steht geschlossen in A250. Dort ist auch verzeichnet, welche Einträge des
Korrekturregisters an welcher Stelle eingearbeitet sind, und welche alten
Dokumente bewusst \emph{keinen} Nachfolger haben, weil ihr Inhalt überholt oder
rein historisch ist.

\section{Sprachstand}

Die A-Serie entsteht zunächst nur auf Deutsch. Das ist eine bewusste Abkehr von
der bisherigen Regel, jede Änderung sofort zweisprachig zu führen. Der Grund
ist arbeitsökonomisch: die Struktur muss sich erst setzen. Eine englische
Fassung entsteht, wenn ein Block abgeschlossen und stabil ist -- dann in einem
Durchgang aus dem fertigen deutschen Stand, nicht Änderung für Änderung.

Bis dahin bleibt die zweisprachige Altserie der Stand für die externe
Korrespondenz.

\section{Zeichenerklärung der A-Serie}

Die folgende Tabelle erklärt die Symbole, die in der gesamten A-Serie durchgehend
verwendet werden. Symbole, die nur in einem einzigen Dokument auftreten, werden
dort beim ersten Auftreten eingeführt.

\begin{center}
\renewcommand{\arraystretch}{1.3}
{\small\setlength\tabcolsep{5pt}
\begin{tabular}{lp{10cm}}
\toprule
\textbf{Symbol} & \textbf{Bedeutung} \\
\midrule
\multicolumn{2}{l}{\textit{Fundamentale Konstanten und Parameter}} \\
$\xi = 4/30000$ & Geometrische Grundkonstante der FFGFT;
  $\xi = 4/3\times10^{-4}$ \\
$\ell_P$ & Planck-Länge, $\ell_P = \sqrt{\hbar G/c^3}
  \approx 1{,}616\times10^{-35}\,$m \\
$\hbar$ & Reduziertes Plancksches Wirkungsquantum \\
$c$ & Lichtgeschwindigkeit (in natürlichen Einheiten: $c=1$) \\
$\alpha$ & Feinstrukturkonstante, $\alpha\approx1/137{,}036$ \\
$\varphi$ & Goldener Schnitt, $\varphi=(1+\sqrt5)/2\approx1{,}618$ \\
\midrule
\multicolumn{2}{l}{\textit{FFGFT-spezifische Größen}} \\
$D_f = 3-\xi$ & Fraktale Dimension des FFGFT-Torus \\
$L_0 = \xi\cdot\ell_P$ & Minimale Längenskala (Sub-Planck),
  $L_0\approx2{,}15\times10^{-39}\,$m \\
$L_T = \ell_P/\xi$ & Torus-Grundperiode, $L_T\approx7500\,\ell_P$ \\
$\tau_0 = \xi\cdot t_P$ & Minimale Zeitskala; per $c=1$ dieselbe Größe wie $L_0$
  in Zeit-Kleidung (Vier-Kleidungen-Prinzip, A080) \\
$K_\text{frak} = 1-100\xi$ & Fraktaler Korrekturfaktor,
  $K_\text{frak}\approx0{,}9867$ \\
$E_0 = \sqrt{m_e m_\mu}$ & Charakteristische Energie,
  $E_0\approx7{,}398\,$MeV \\
$\tilde T\cdot m = 1$ & Zeit-Masse-Dualität (Setzung); $\tilde T$ ist das
  intrinsische Zeitfeld \\
$v = 246\,$GeV & Higgs-Vakuumerwartungswert (Energieeinheit) \\
\midrule
\multicolumn{2}{l}{\textit{Schichtmarkierungen}} \\
\textbf{[SETZUNG]} & Axiom der Theorie; wird nicht weiter begründet \\
\textbf{[B]} & Brücke; algebraisch bewiesen \\
\textbf{[K]} & Kern; numerisch geprüfte Ableitung \\
\textbf{[S]} & Skizze; plausibel, aber nicht vollständig ausgeführt \\
\midrule
\multicolumn{2}{l}{\textit{Mathematische Strukturen}} \\
$T^4/\mathbb{Z}_3$ & Kompakter 4-Torus mit $\mathbb{Z}_3$-Orbifold-Struktur \\
$\mathcal{H} = L^2(T^4)\otimes\mathbb{C}^3$ & Hilbertraum der Theorie \\
$\pi_1(T^4)\cong\mathbb{Z}^4$ & Erste Homotopiegruppe des 4-Torus \\
$P_i = |\mathbf{e}_i\rangle\langle\mathbf{e}_i|$ & Projektor auf Mode $i$ \\
$\Phi = \sum_i m_i P_i$ & Moden-Operator (Feldoperator) \\
\midrule
\multicolumn{2}{l}{\textit{Massen und Verhältnisse}} \\
$m_e, m_\mu, m_\tau$ & Elektron-, Myon-, Tauon-Masse \\
$m_i = r_i\cdot\xi^{p_i}\cdot v$ & FFGFT-Massenformel; $r_i\in\mathbb{Q}$
  rationaler Vorfaktor, $p_i\in\mathbb{Q}$ Exponent \\
$(n, l, j)$ & Quantenzahlen: Haupt-, Bahn-, Spinquantenzahl \\
\bottomrule
\end{tabular}
}
\renewcommand{\arraystretch}{1.0}
\end{center}

\section{Eigenständige Dokumente des Altbestands}

\textbf{[SETZUNG]} Die A-Serie übernimmt 102 Dokumente des Altbestands. Der
Rest -- mehr als 200 Nummern -- bleibt eigenständig. Dieser Abschnitt listet
die wichtigsten thematischen Gruppen dieser eigenständigen Dokumente, damit der
Leser weiß, wo der Altbestand endet und wo die A-Serie anfängt. Vollständige
Zuordnung: A250.

\textbf{Kosmologischer Sektor (ausgeschlossen, Dok.~036, 063, 064, 140 u.a.).}
Hubble-Konstante, kosmische Mikrowellenhintergrundstrahlung, Vakuumfeld auf
galaktischen Skalen, $\Lambda_\text{cosmo}$. Dieser Sektor ist nicht falsch --
er ist von anderem Anspruch: er arbeitet durchgehend mit importierten Skalen und
Kalibrierungen, die die A-Serie nicht reproduzieren kann. Ausgeschlossen, nicht
überholt.

\textbf{Eichsektor und QFT-Erweiterungen (Dok.~186, 192, 281 u.a.).}
Vollständige Feynman-Regeln, Renormierungsgruppe jenseits der $\xi$-Korrekturen,
Yang-Mills-Formulierung. Diese Dokumente behandeln Themen, die in der A-Serie als
offene Punkte deklariert sind (A190, A240). Sie bleiben eigenständig, weil die
Projektionsabbildung noch fehlt.

\textbf{IPI-Rahmenvergleiche (Dok.~246, 251, 252, 265, 267, 269, 271, 272,
274, 283, 285, 294, 297, 299 u.a.).} Strukturvergleiche mit HLV (Helix-Light-Vortex),
RA/PMT, Dot Theory, SICC und anderen Arbeiten aus dem IPI-Umfeld. Diese Dokumente
sind explizit keine FFGFT-Aussagen, sondern Vergleichsdokumente: sie untersuchen,
wo externe Rahmen und FFGFT Berührungspunkte haben. Die A-Serie übernimmt keine
dieser Vergleiche (A240 nennt die Berührungspunkte, ohne sie zu tragen).

\textbf{Quanten-Hardware und photonische Systeme (Dok.~083--085, 161, 167,
173--176, 183 u.a.).} Photonische Chips in der Nachrichtentechnik, Quantum Dots,
Shors Algorithmus auf photonischer Plattform, Google-Willow-Analyse, Ising-Maschinen.
Diese Dokumente wenden FFGFT-Ideen auf konkrete Hardwaresysteme an. Die Anwendung
ist eigenständig und nicht rückwärts-verankert in den A-Serien-Kern.

\textbf{Populäre und narrative Einführungen (Dok.~002, 004, 185, 191 u.a.).}
\enquote{Die Reise} (narrative Einführung ohne Formeln), FFGFT für Kommunikationsingenieure,
Videoscripte. Diese Dokumente sind bewusst nicht Teil der A-Serie -- sie sind
der Zugang für Leser ohne Physik-Vorkenntnisse und bleiben es. Sie werden nicht
überholt, sondern ergänzt.

\textbf{KI und Bewusstsein (Dok.~100, 207 u.a.).} T0-Theorie und Bewusstsein;
Bewusstsein und Brücken. Diese Dokumente behandeln ein Thema, das in der A-Serie
nicht vorkommt. Sie sind eigenständige philosophische Arbeiten, kein Anhang der Physik.

\textbf{Methodische und diagnostische Dokumente (Dok.~241, 242, 244, 247, 248,
250, 262, 277, 281, 284, 287, 288, 303--305 u.a.).} Hawking-Informationsparadoxon,
Verhältnis- vs. Absolutwert-Pädagogik, Dokumentenstrukturkarten, Einbettungs-Protokolle.
Diese Dokumente sind Arbeitsunterlagen, die während der Entwicklung des Korpus
entstanden. Sie sind für den Nachvollzug der Entwicklungsgeschichte wichtig,
aber kein Teil des kanonischen Stands.

\textbf{Physik-Erweiterungen im engeren Sinn (Dok.~021--024, 027--034,
037--040, 054--068, 075--081, 089, 095, 105, 114, 124, 129, 137, 145--170,
172, 178, 179, 182, 184, 187, 192, 206 u.a.).} Ältere Formulierungen der Massenformel,
alternative Parametrisierungen, Zwischenstadien der Theorie, Einzelrechnungen die
in A-Dokumenten aufgegangen sind oder bewusst nicht übernommen wurden. Dieser
Block ist am größten und am heterogensten: er enthält frühe korrekte Rechnungen
neben überholten Ansätzen, eigenständige Nebenlinien neben bereits aufgenommenen
Substanzen. Das Korrekturregister (Dok.~190) bucht die relevanten Korrekturen;
die A-Dokumente nennen, was sie übernehmen. Was weder im Register steht noch
in einem A-Dokument genannt wird, ist historischer Bestand.


\section{Was offen bleibt}

Dieser Abschnitt steht in jedem Dokument der A-Serie und benennt, was das
Dokument \emph{nicht} leistet. Die Regel ist von außen übernommen: ein
abgeschlossener Schritt, der sein Residuum nicht benennt, beansprucht
stillschweigend Zuständigkeit über das hinaus, was er gezeigt hat.

Für dieses Dokument sind vier Punkte offen.

\textbf{Erstens deckt die Serie nicht den ganzen Korpus ab.} Sie ordnet den
geometrischen Kern, die Einheiten- und Konstantenstruktur und den Fermionsektor;
der Eichsektor und die dynamische Feldtheorie bleiben im Altbestand, ebenso der
aus der Serie ausgeschlossene Sektor der frühen Kosmologie (A240, A250). Die Zusage, jede rechnerische Aussage mit einem
Prüfskript zu
hinterlegen, ist für die Serie als Ganzes noch nicht eingelöst.

\textbf{Zweitens ist die Ersetzungsbehauptung ungeprüft.} Der Bauplan ordnet
jedem A-Dokument eine Menge alter Nummern zu, deren Inhalt es aufnehmen soll.
Ob diese Zuordnung vollständig ist -- ob also nichts aus dem Altbestand
herausfällt --, lässt sich erst am fertigen Register A250 entscheiden, nicht
vorher. Bis dahin ist \enquote{ersetzt} eine Absicht, keine Feststellung.

\textbf{Drittens ist die Grenze des Kosmologie-Ausschlusses nicht scharf.} Der
Ausschluss ist als Trennung nach Anspruchsniveau begründet, nicht nach Thema.
An drei Stellen -- der Korrekturfaktor $K_\text{frak}$, die Brückenkonstanten,
die Projektionskette -- trägt der Altbestand historisch kosmologische Anker.
Ob diese Stellen ohne solche Anker tragfähig bleiben, ist eine offene
inhaltliche Frage und nicht durch die Ausschlussentscheidung beantwortet.

\textbf{Viertens ist der Toleranzmaßstab hier nur formuliert, nicht
angewandt.} Die Forderung, jede Zahlenaussage gegen Messgenauigkeit und
Bestimmungstoleranz zu stellen, ist leicht auszusprechen. Ob sie durchhaltbar
ist, entscheidet sich erst dort, wo Zahlen stehen -- also ab A100.

\section{Ersetzt}

Das genannte Korrekturregister ist Dok.~190 (Einträge K1--K6, P1--P43, R1--R56); die Wicklungszahl-Korrekturen stehen in Dok.~210.

Dieses Dokument nimmt die Rahmen- und Übersichtsteile folgender Dokumente auf:
Dok.~086 (Dokumentenübersicht), Dok.~201 (FFGFT-alles, Gliederungsteil),
Dok.~205 (Narrativ, Rahmenabschnitte). Der narrative Teil von Dok.~205 bleibt
eigenständig lesenswert und ist nicht ersetzt.
